\begin{abstract}
Lung adenocarcinoma (LUAD) exhibits substantial inter-patient molecular heterogeneity
that transcriptomic subtyping of tumour tissue only partially captures.
Here we ask whether this heterogeneity extends to the tumourigenesis \emph{process}
itself: do patients undergo qualitatively distinct molecular transformations?
We apply Multi-Omics Factor Analysis (MOFA+) to 102 matched tumour-normal pairs from the
Clinical Proteomic Tumour Analysis Consortium LUAD cohort, learning a 13-factor latent
space from transcriptomic, proteomic, and phosphoproteomic data.
For each patient we compute a latent transformation vector
$\Delta Z = Z_{\text{tumour}} - Z_{\text{normal}}$,
encoding the direction and magnitude of molecular change in the shared latent space.
Clustering these vectors by k-means identified three transformation subtypes that are
reproducible across alternative representations (ARI\,$\geq$\,0.80).
One-vs-rest differential expression followed by preranked GSEA revealed systematic
differences in transcriptional reprogramming: subtype~0 shows relatively suppressed
proliferative reprogramming and is significantly enriched for \textit{EGFR} mutations;
subtype~1 is dominated by E2F and G2-M cell-cycle activation; and subtype~2 combines
interferon signalling with a proliferative component.
These signatures are broadly consistent with the TRU, PP, and PI state-based LUAD
subtypes respectively, though the correspondence is suggestive rather than formally
established.
No significant overall survival difference was detected across subtypes ($n = 102$).
The analysis provides a proof-of-concept that latent transformation vectors derived
from a matched multi-omics cohort can reveal biologically coherent modes of
tumourigenesis.
\end{abstract}
